\section{String, Kaluza--Klein, and \texorpdfstring{$G_2$}{G2} completion problem}
\label{sec:string-kk-g2}

The secular equation has the form of a low-level branch splitting that preserves a high-spin datum.  This is the reason string and Regge language enter the project.  In string theory, higher-spin excitations organize into Regge trajectories, and the asymptotic slope is controlled by the string tension.  The Rivero construction was motivated in part by preserving such an asymptotic datum while modifying low-spin masses \cite{Rivero2006,TongString}.

A Kaluza--Klein or brane derivation would have to produce the determinant
\begin{equation}
  \det\begin{pmatrix}
    x & -\sqrt{J}\\
    -\sqrt{J} & x+J
  \end{pmatrix}=0.
  \label{eq:det-completion}
\end{equation}
This can arise from a two-channel boundary problem, a mixing of brane-localized and bulk modes, or an effective orbit problem.  The required output is not merely a qualitative similarity; it is the exact \(J\)-dependence in Eq.~\eqref{eq:det-completion}.

Extra-dimensional electroweak models show that boundary conditions can break gauge symmetry and determine vector spectra.  Higgsless constructions on intervals provide a controlled language for relating boundary data, vector masses, fermion localization, and precision observables \cite{CsakiHubiszMeade2005}.  In the present project, such models serve as templates for deriving a mass-ratio constraint from boundary conditions.

M-theory on \(G_2\)-holonomy manifolds supplies a separate completion map.  Smooth compact \(G_2\) compactifications give four-dimensional minimal supersymmetry but do not by themselves yield realistic charged chiral matter in the large smooth limit.  Singularities are required for nonabelian gauge sectors and chiral fermions \cite{WittenG2Anomaly2001,AcharyaWitten2001,AcharyaGukov2004}.  Local Higgs-bundle descriptions over associative three-cycles give a modern way to organize gauge sectors, localized matter, and singular transitions \cite{BraunEtAl2019}.

The completion problem can therefore be stated as a table of mechanisms.
\begin{center}
\begin{tabular}{lll}
\toprule
Mechanism & What it can supply & Required DeVries test\\
\midrule
Regge/string tower & high-spin asymptotics & preserve slope while producing Eq.~\eqref{eq:det-completion}\\
Interval/brane BC & vector mass spectra & derive \(J=3/4,2\) samples from boundary data\\
KK compactification & mode eigenvalues & produce two-channel branch determinant\\
Singular \(G_2\) & gauge/chiral sectors & localize the electroweak operator geometrically\\
Higgs bundle & matter/mass localization & identify the two branches as spectral-cover data\\
\bottomrule
\end{tabular}
\end{center}

\paragraph*{Section status.}  This section states completion mechanisms and tests.  No dynamical derivation is claimed here.
